\chapter{分数初步}\label{ch:g6:fractions}

把三个披萨公平地分给四个人，每人得到\,\dots\
不是整数个披萨。\omterm{def:g6:fractions:def}{分数}正是为“分配”而发明的数。
本章建立\omterm{def:g6:fractions:def}{分数}的图像：\omterm{def:g6:fractions:def}{分数}作为整体的一部分，
作为数轴上的一个点，以及
除法的精确结果。

\section{分数的含义}

\begin{definition}[分数]\label{def:g6:fractions:def}
将一个单位平均分成 $b$ 份，取其中 $a$ 份，所得的量
即为\emph{分数}\index{分数}：
\[
\frac{a}{b}
\qquad
\begin{array}{l}
a \text{：\emph{分子}——取走的份数；}\\
b \text{：\emph{分母}——单位被平均分成的份数。}
\end{array}
\]
\end{definition}

\begin{omfigure}
\begin{tikzpicture}[scale=0.55]
  \foreach \i in {0,...,3} \draw[omExo] (\i*2,0) rectangle (\i*2+2,1.2);
  \foreach \i in {0,...,2} \fill[omDef!40] (\i*2+0.05,0.05)
    rectangle (\i*2+1.95,1.15);
  \draw[thick] (0,0) rectangle (8,1.2);
  \node[below, font=\small] at (4,-0.15) {$\frac34$ 的条形：平均分成
  $4$ 份，取 $3$ 份};
  \begin{scope}[xshift=11.5cm]
  \foreach \i in {0,...,7} \draw[omExo] (\i,0) rectangle (\i+1,1.2);
  \foreach \i in {0,...,4} \fill[omThm!35] (\i+0.05,0.05)
    rectangle (\i+0.95,1.15);
  \draw[thick] (0,0) rectangle (8,1.2);
  \node[below, font=\small] at (4,-0.15) {$\frac58$ 的条形：平均分成
  $8$ 份，取 $5$ 份};
  \end{scope}
\end{tikzpicture}

{\small 同一条形的两个\omterm{def:g6:fractions:def}{分数}。\omterm{def:g4:fractions:def}{分母}说明切分得有多细，
\omterm{def:g4:fractions:def}{分子}说明取了多少。}
\end{omfigure}

\begin{example}\label{ex:g6:fractions:morethanone}
\omterm{def:g6:fractions:def}{分数}也可以大于 $1$：$\frac74$ 表示七个\omterm{def:g3:division:half}{四分之一}——
即一个整单位（四个\omterm{def:g3:division:half}{四分之一}）加上多出的三个\omterm{def:g3:division:half}{四分之一}：
\[
\frac74 = 1 + \frac34 .
\]
\end{example}

\begin{proposition}[数轴上的分数]\label{prop:g6:fractions:line}
要在数轴上标出 $\frac ab$，只需将每个单位区间分成 $b$
等份，数出 $a$ 份，从 $0$ 算起。
\end{proposition}

\begin{omfigure}
\begin{tikzpicture}[scale=1.5]
  \draw[-Stealth] (-0.2,0) -- (7.6,0);
  \foreach \i in {0,4,8,12,16,20,24,28}
    \draw ({\i*0.25},-0.09) -- ({\i*0.25},0.09);
  \foreach \i in {0,...,28} \draw ({\i*0.25},-0.05) -- ({\i*0.25},0.05);
  \node[below, font=\small] at (0,-0.1) {$0$};
  \node[below, font=\small] at (1,-0.1) {$1$};
  \node[below, font=\small] at (2,-0.1) {$2$};
  \node[below, font=\small] at (3,-0.1) {$3$};
  \node[below, font=\small] at (4,-0.1) {$4$};
  \node[below, font=\small] at (5,-0.1) {$5$};
  \node[below, font=\small] at (6,-0.1) {$6$};
  \node[below, font=\small] at (7,-0.1) {$7$};
  \fill[omDef] (0.75,0) circle (1.6pt)
    node[above, font=\small] {$\frac34$};
  \fill[omThm] (1.75,0) circle (1.6pt)
    node[above, font=\small] {$\frac74$};
  \fill[omMeth] (5.5,0) circle (1.6pt)
    node[above, font=\small] {$\frac{22}{4}$};
\end{tikzpicture}

{\small 数轴按\omterm{def:g3:division:half}{四分之一}刻度：$\frac34$ 位于 $1$ 之前，
$\frac74$ 位于 $1$ 与 $2$ 之间，$\frac{22}{4} = 5 + \frac24$ 位于
$5$ 与 $6$ 的正中间。}
\end{omfigure}

\begin{theorem}[分数即除法]\label{thm:g6:fractions:quotient}
\omterm{def:g6:fractions:def}{分数} $\frac ab$ 是这样一个数：乘以 $b$ 后
得到 $a$：
\[
\frac ab \times b = a .
\]
换句话说，$\frac ab$ 是 $a$ 除以 $b$ 的精确商——即使
除法“除不尽”时也是如此。
\end{theorem}
\admitted

\begin{example}\label{ex:g6:fractions:pizzas}
三个披萨分给四个人：每人得到 $\frac34$ 个披萨。
用定理验证：$4$ 份 $\frac34$ 合起来为
$\frac34 \times 4 = 3$ 个披萨。有些\omterm{def:g6:fractions:def}{分数}是\omterm{def:g6:decimals:places}{小数}
（$\frac34 = 0.75$），另一些则不是：$\frac13 = 0.333\dots$ 永远不会结束。
\omterm{def:g6:fractions:def}{分数}是\emph{精确}的值，$0.33$ 只是近似值。
\end{example}

\section{等值分数}

\begin{proposition}[等值分数]\label{prop:g6:fractions:equal}
当\omterm{def:g6:fractions:def}{分数}的\omterm{def:g4:fractions:def}{分子}与\omterm{def:g4:fractions:def}{分母}同时
乘以（或除以）同一个非零数时，\omterm{def:g6:fractions:def}{分数}的值不变：
\[
\frac{a}{b} = \frac{a \times k}{b \times k} .
\]
\end{proposition}

\begin{proof}[证明思路]
将分得的每一份再切成 $k$ 个更小的部分，
使所取的份数与总份数都乘以 $k$，
但所取的量保持不变。图示最能说明问题：
\end{proof}

\begin{omfigure}
\begin{tikzpicture}[scale=0.6]
  \foreach \i in {0,1,2} \draw[omExo] (\i*2.4,0) rectangle (\i*2.4+2.4,1.2);
  \foreach \i in {0,1} \fill[omDef!40] (\i*2.4+0.05,0.05)
    rectangle (\i*2.4+2.35,1.15);
  \draw[thick] (0,0) rectangle (7.2,1.2);
  \node[below, font=\small] at (3.6,-0.15) {$\frac23$};
  \begin{scope}[xshift=10.5cm]
  \foreach \i in {0,...,5} \draw[omExo] (\i*1.2,0) rectangle (\i*1.2+1.2,1.2);
  \foreach \i in {0,...,3} \fill[omDef!40] (\i*1.2+0.05,0.05)
    rectangle (\i*1.2+1.15,1.15);
  \draw[thick] (0,0) rectangle (7.2,1.2);
  \node[below, font=\small] at (3.6,-0.15) {$\frac46$};
  \end{scope}
\end{tikzpicture}

{\small $\frac23 = \frac46$：将每个三分之一平分为两份，使\omterm{def:g4:fractions:def}{分子}、\omterm{def:g4:fractions:def}{分母}都翻倍，
阴影部分的大小不变。}
\end{omfigure}

\begin{method}[约分]\label{met:g6:fractions:simplify}
寻找一个数，它能\emph{同时}\omterm{def:g6:wholes:divisible}{整除}\omterm{def:g4:fractions:def}{分子}与\omterm{def:g4:fractions:def}{分母}，
并将两者都除以这个数。重复此过程，直到不再有公\omterm{def:g5:division:multiple}{因数}为止：
\[
\frac{12}{18} = \frac{12 \div 2}{18 \div 2} = \frac{6}{9}
= \frac{6 \div 3}{9 \div 3} = \frac{2}{3}.
\]
\end{method}

\section{求一个量的几分之几}

\begin{method}[求一个量的几分之几]\label{met:g6:fractions:ofquantity}
求一个量的 $\frac ab$ 时，先将该量除以 $b$（得到一份的大小），
再乘以 $a$（即取的份数）：
\[
\frac ab \text{ 的 } Q = (Q \div b) \times a .
\]
（先乘以 $a$ 再除以 $b$，所得结果相同，
选择便于计算的顺序即可。）
\end{method}

\begin{example}\label{ex:g6:fractions:ofquantity}
计算 $\frac35$ 的 $40$ 欧元是多少，分步进行：
\begin{enumerate}
  \item $40$ 的五分之一：$40 \div 5 = 8$；
  \item 五分之三：$8 \times 3 = 24$。
\end{enumerate}
因此 $\frac35$ 的 $40$ 欧元是 $24$ 欧元。验算：剩下的
$\frac25$ 为 $16$ 欧元，而 $24 + 16 = 40$。
\end{example}

\section{练习}

\begin{exercise}[$\star$]\label{exo:g6:fractions:1}
根据下列图形描述，写出对应的\omterm{def:g6:fractions:def}{分数}：一块蛋糕切成 $8$ 块，
吃掉 $3$ 块；一块巧克力共 $12$ 格，吃掉 $7$ 格；一个派
切成 $6$ 块，$6$ 块全部吃掉。
\end{exercise}

\begin{exercise}[$\star$]\label{exo:g6:fractions:2}
画一条长条，平均分成 $5$ 份，并涂出其中的 $\frac45$，
判断 $\frac45$ 比 $1$ 小还是大？$\frac65$ 呢？
\end{exercise}

\begin{exercise}[$\star$]\label{exo:g6:fractions:3}
在按三分之一刻度的数轴上标出：$\frac13$；\ $\frac53$；
\ $2$；\ $\frac73$。这些数中，哪两个到
$2$ 的距离相等？
\end{exercise}

\begin{exercise}[$\star$]\label{exo:g6:fractions:4}
将每个\omterm{def:g6:fractions:def}{分数}写成一个整数与一个小于 $1$ 的\omterm{def:g6:fractions:def}{分数}的和
（如 $\frac74 = 1 + \frac34$）：
\[
\frac{9}{4}, \qquad \frac{13}{5}, \qquad \frac{12}{6}, \qquad
\frac{25}{8} .
\]
\end{exercise}

\begin{exercise}[$\star$]\label{exo:g6:fractions:5}
填空，使下列\omterm{def:g6:fractions:def}{分数}相等：
\[
\frac{2}{3} = \frac{?}{12}, \qquad
\frac{15}{20} = \frac{3}{?}, \qquad
\frac{7}{10} = \frac{21}{?}, \qquad
\frac{?}{6} = \frac{10}{12}.
\]
\end{exercise}

\begin{exercise}[$\star$]\label{exo:g6:fractions:6}
将下列\omterm{def:g6:fractions:def}{分数}化为最简形式：$\dfrac{6}{8}$；\quad $\dfrac{15}{25}$；
\quad $\dfrac{18}{24}$；\quad $\dfrac{35}{7}$。
\end{exercise}

\begin{exercise}[$\star$]\label{exo:g6:fractions:7}
计算：$\frac12$ 的 $86$；\quad $\frac34$ 的 $60$；\quad $\frac25$ 的
$35$；\quad $\frac{7}{10}$ 的 $250$。
\end{exercise}

\begin{exercise}[$\star$]\label{exo:g6:fractions:8}
哪个更大：$\frac12$ 还是 $\frac13$（同一块蛋糕的）？$\frac25$
还是 $\frac35$？$\frac34$ 还是 $\frac78$（将两者都化为\omterm{def:g4:fractions:def}{分母}为
$8$ 的\omterm{def:g6:fractions:def}{分数}）？
\end{exercise}

\begin{exercise}[$\star\star$]\label{exo:g6:fractions:9}
  写出能化为\omterm{def:g6:decimals:places}{小数}的\omterm{def:g6:fractions:def}{分数}的\omterm{def:g6:decimals:places}{小数}值，并说明哪个\omterm{def:g6:fractions:def}{分数}不能化为\omterm{def:g6:decimals:places}{小数}：
$\frac12$；\ $\frac34$；\ $\frac13$；\ $\frac{7}{10}$；
\ $\frac{9}{4}$。
\end{exercise}

\begin{exercise}[$\star\star$]\label{exo:g6:fractions:10}
某班有 $28$ 名学生，其中四分之三乘校车上学，其余学生步行。
步行的学生有多少人？请分步求解，
并检验两组人数之和是 $28$。
\end{exercise}

\begin{exercise}[$\star\star$]\label{exo:g6:fractions:11}
利奥花掉了自己积蓄的 $\frac23$ 去买一款售价 $18$ 欧元的游戏。
利奥原来有多少钱？（积蓄的三分之一是 $18$ 的\omterm{def:g3:division:half}{一半}\,\dots\ 想一想为什么，
或者先求出一份的大小。）
\end{exercise}

\begin{exercise}[$\star\star\star$]\label{exo:g6:fractions:12}
一个水箱注水至其\omterm{def:g2:measure:units}{容积}的 $\frac58$。用掉 $10$ 升水后，
水箱恰好半满。水箱的\omterm{def:g2:measure:units}{容积}是多少？
（$10$ 升水占水箱的几分之几？）
\end{exercise}

\section{问题：你永远吃不完的巧克力棒}

\begin{problem}[{周末问题——一半的一半：和
$\frac12 + \frac14 + \frac18 + \dots$ 不断逼近 $1$，
却永远无法触及它}]
\label{pb:g6:fractions:1}
拿一根巧克力棒，吃掉其中的\omterm{def:g3:division:half}{一半}，再吃掉\emph{剩下部分}的\omterm{def:g3:division:half}{一半}，
接着再吃掉剩下部分的\omterm{def:g3:division:half}{一半}，如此继续下去。看上去有两件事同时成立：
你永远吃不完这根巧克力棒（总有一些剩下），但你却几乎吃掉了全部。
本题把这两种感受化为关于\omterm{def:g6:fractions:def}{分数}的精确论断，途中还会遇到
一位古希腊的赛跑者，以及一道
完全用一把只会把东西对半切开的小刀
来解决的分配难题。

\medskip
\noindent\textbf{第一部分——\omterm{def:g3:division:half}{一半}的\omterm{def:g3:division:half}{一半}。}
这根巧克力棒是由 $16$ 个相等的正方形组成的长方形。
\begin{enumerate}
  \item 第一口：吃掉\omterm{def:g3:division:half}{一半}的巧克力棒，是多少个正方形？
        用两种方式把这一口写成占整根巧克力棒的\omterm{def:g6:fractions:def}{分数}，
        分别以\omterm{def:g4:fractions:def}{分母} $2$ 和\omterm{def:g4:fractions:def}{分母} $16$ 表示，
        （见\cref{prop:g6:fractions:equal}）。
  \item 第二口：吃掉剩下部分的\omterm{def:g3:division:half}{一半}，是多少个正方形？
        借助图形或等值\omterm{def:g6:fractions:def}{分数}，说明这一口是整根巧克力棒的 $\frac14$：
        \emph{\omterm{def:g3:division:half}{一半}的\omterm{def:g3:division:half}{一半}
        就是\omterm{def:g3:division:half}{四分之一}}。
  \item 第三口和第四口规则相同：分别写出各自的正方形个数，
        以及占整根巧克力棒的\omterm{def:g6:fractions:def}{分数}。
        相邻两口之间，\omterm{def:g4:fractions:def}{分母}有什么变化？
  \item 吃完四口后，一共吃掉了多少个正方形？
        这些占整根巧克力棒的几分之几？
        还剩下几分之几？
  \item 第 4 问的结论用十六分之一来写，就是：
        \[
        \frac{8}{16} + \frac{4}{16} + \frac{2}{16} + \frac{1}{16}
        = \frac{15}{16},
        \qquad\text{即}\qquad
        \frac12 + \frac14 + \frac18 + \frac{1}{16}
        = \frac{15}{16} .
        \]
        请逐一检验这四个改写式
        （$\frac12 = \frac{8}{16}$，依此类推）。
\end{enumerate}

\medskip
\noindent\textbf{第二部分——通向 $1$ 的阶梯。}
为了继续吃下去，设想一根由 $64$ 个相等正方形组成的更精细的巧克力棒。
\begin{enumerate}[resume]
  \item 列出前六口的大小（以正方形为单位），
        依次为（$32$，然后\,\dots），并检验第六口之后
        恰好只剩下一个正方形。剩下的部分占整根
        巧克力棒的几分之几？
  \item 每一口之后，巧克力棒还剩下几分之几？写出
        直到第六口为止的列表（$\frac12$，$\frac14$，\dots），
        并描述其中的规律。若想象中的一把刀
        甚至能把最后一个正方形也对半切开，
        那么第七口之后还会剩下多少？
  \item 解释为什么无论吃多少口，吃掉的总量\emph{永远}
        不可能达到 $1$。（与\cref{pb:g6:decimals:1}中的
        二十个九相比，是否是同一个故事？）
  \item 然而，吃掉的总量会超过任何小于 $1$ 的目标。一位朋友
        向你挑战，要吃掉超过 $\frac{99}{100}$ 的巧克力棒。
        请比较 $\frac{1}{128}$ 与
        $\frac{1}{100}$ 哪个更小，
        并说明在第几口之后挑战成功？
  \item 在一条从 $0$ 到 $1$、以十六分之一为刻度的数轴上，
        标出前四口各自之后吃掉的总量：
        $\frac12$，$\frac34$，$\frac78$，$\frac{15}{16}$。
        每一个新的总量落在何处，与前一个总量相比，
        以及与 $1$ 相比，各在什么位置？
\end{enumerate}

\medskip
\noindent\textbf{第三部分——芝诺的赛跑者，
以及只会对半切的刀。}
\begin{enumerate}[resume]
  \item 二十四个世纪前，希腊哲学家芝诺讲述过这样一个故事：要走到
        一堵墙边，赛跑者必须先跑完\omterm{def:g3:division:half}{一半}的距离，再跑完剩下距离的\omterm{def:g3:division:half}{一半}，
        再跑完剩下距离的\omterm{def:g3:division:half}{一半}，如此继续下去\,\dots\
        “所以赛跑者永远到不了墙边！”利用第 8 题和第 9 题，说明芝诺故事中
        哪些是\emph{正确}的，哪里是陷阱。
        （赛跑者真的需要
        无限长的时间
        才能完成所有这些步骤吗？）
  \item 分配难题：把 $15$ 根相同的巧克力棒公平地分给 $16$ 个孩子，
        只允许使用对半切割，
        （对整根、半根、\omterm{def:g3:division:half}{四分之一}根\,\dots\ 进行切割）。
        请说明如何切割：有多少根切成两半，
        有多少根切成\omterm{def:g3:division:half}{四分之一}，有多少根切成八分之一、十六分之一，
        使得每个孩子都能得到每种大小的一块。
  \item 检验你的分配：把每个孩子的份额写成若干\omterm{def:g6:fractions:def}{分数}之和，
        以十六分之一为单位计算，
        并验证这 $16$ 份合起来恰好用掉 $15$ 根巧克力棒。
  \item 隔壁的班级里，$7$ 根巧克力棒被公平地分给 $8$ 个孩子。
        哪个孩子得到的巧克力更多——是 $16$ 份中的一份，
        还是 $8$ 份中的一份？（将两份都化为\omterm{def:g4:fractions:def}{分母}为 $16$ 的\omterm{def:g6:fractions:def}{分数}，
        见\cref{exo:g6:fractions:8}。）
  \item 最后一问：一位朋友声称，若咬下\emph{无穷多}口，
        就会\emph{恰好}吃掉整根巧克力棒。画一个正方形，
        在其中涂出各口：先涂去正方形的\omterm{def:g3:division:half}{一半}，
        再涂去\omterm{def:g3:division:half}{四分之一}，然后是八分之一\,\dots\ 你观察到
        未涂色的角落有什么特点？请仔细陈述本题中的
        两个正确的事实：每吃掉有限多口之后还剩下多少，
        以及这一剩余部分会变得多么小。
        （“无穷和”究竟是什么意思，
        要留待高中分册中再讲。）
\end{enumerate}
\end{problem}
